% Part: first-order-logic
% Chapter: sequent-calculus
% Section: equality

\documentclass[../../../include/open-logic-section]{subfiles}

\begin{document}

\olfileid{fol}{seq}{ide}

\olsection{\usetoken{P}{derivation} met \usetoken{S}{identity}}

Voor !!^{derivation}s waarin !!{identity}, dus gelijkheid, voorkomt,
zijn extra beginsequenten en afleidingsregels nodig.

\begin{defn}[Beginsequenten voor $\eq$]
Als $t$ een gesloten term is, dan is ${} \Sequent \eq[t][t]$ een
beginsequent. Deze sequent drukt
zelfidentiteit uit.
\end{defn}

De regels voor $\eq$ staan hieronder. Daarbij zijn $t_1$ en $t_2$
gesloten termen.

\begin{defish}
\Axiom$ \eq[t_1][t_2], \Gamma \fCenter \Delta, !A(t_1) $
\RightLabel{$\eq$}
\UnaryInf$\eq[t_1][t_2], \Gamma \fCenter \Delta, !A(t_2)$
\DisplayProof
\hfill
\Axiom$\eq[t_1][t_2], \Gamma \fCenter \Delta, !A(t_2) $
\RightLabel{$\eq$}
\UnaryInf$\eq[t_1][t_2], \Gamma  \fCenter \Delta, !A(t_1)$
\DisplayProof
\end{defish}

\begin{ex}
Als $s$ en $t$ gesloten termen zijn, dan geldt $\eq[s][t], !A(s)
\Proves !A(t)$:
\begin{prooftree}
\Axiom$ !A(s) \fCenter !A(s)$
\RightLabel{\LeftR{\Weakening}}
\UnaryInf$\eq[s][t], !A(s)  \fCenter !A(s)$
\RightLabel{$\eq$}
\UnaryInf$\eq[s][t], !A(s)  \fCenter !A(t)$
\end{prooftree}
Dit heet het beginsel van substitueerbaarheid van identieken, of de
wet van Leibniz: gelijke termen mogen op deze manier voor elkaar
worden vervangen.

$\Log{LK}$ bewijst ook twee eigenschappen van $\eq$: gelijkheid is
symmetrisch en transitief.
\begin{prooftree}
\Axiom$ \fCenter \eq[t_1][t_1] $
\RightLabel{\LeftR{\Weakening}}
\UnaryInf$ \eq[t_1][t_2] \fCenter \eq[t_1][t_1] $
\RightLabel{$\eq$}
\UnaryInf$ \eq[t_1][t_2] \fCenter \eq[t_2][t_1]$
\DisplayProof\qquad\bottomAlignProof
\Axiom$ \eq[t_1][t_2] \fCenter \eq[t_1][t_2] $
\RightLabel{\LeftR{\Weakening}}
\UnaryInf$\eq[t_2][t_3], \eq[t_1][t_2]  \fCenter \eq[t_1][t_2] $
\RightLabel{$\eq$}
\UnaryInf$\eq[t_2][t_3], \eq[t_1][t_2]  \fCenter \eq[t_1][t_3]$
\RightLabel{\LeftR{\Exchange}}
\UnaryInf$\eq[t_1][t_2], \eq[t_2][t_3]  \fCenter \eq[t_1][t_3]$
\end{prooftree}
Bij de linker !!{derivation} vullen we voor de
!!{formula}~$!A(x)$ de formule~$\eq[x][t_1]$ in. Bij de rechter vullen
we voor $!A(x)$ de formule~$\eq[t_1][x]$ in. Zo gebruikt elk bewijs
dezelfde gelijkheidsregel
in een andere richting.
\end{ex}

\begin{prob}
Werk voor elk van de volgende sequenten een !!{derivation} uit:
\begin{enumerate}
\item $\Sequent \lforall[x][\lforall[y][((x = y \land !A(x)) \lif !A(y))]]$
\item $\lexists[x][!A(x)] \land \lforall[y][\lforall[z][((!A(y) \land
    !A(z)) \lif y = z)]] \Sequent 
\lexists[x][(!A(x) \land \lforall[y][(!A(y) \lif y = x)])]$
\end{enumerate}
\end{prob}
\end{document}
